\chapter{INTRODUCTION}

\section{Background}
\setlength{\parindent}{0pt}
\hspace{1.6 cm}Finding an available time slot between friends is also a time-consuming struggle as students need to look up each other’s schedule constantly. Academic planning is especially important for engineering students who must follow a structured curriculum with multiple prerequisites and limited courses offered each semester. Evaluating course difficulty, reviewing the course curriculum and coordinating schedules with peers for group study are also the struggles that most students are facing in the faculty. At the EGCI program, students currently lack a centralized platform that allows them to find available time through class-schedule screenshot, plan their academic paths in a visual and structured way, and  rate the course. 

\hspace{1.6 cm}Comparing schedules with friends to find common free time requires manual checking which is time consuming and error prone when there are multiple students. The process of just using the screenshots of the schedule and finding the available times would be a great solution for  meetings and collaborations.  In addition, the issues of the manual schedule coordinations of students can also be solved by the combination of TableExtractNet, by Ngubane and Tapamo \cite{ngubane2024tableextractnet} a deep learning approach utilizing CornerNet and Faster R-CNN to detect and recognize table structures from document images, which can be implemented and trained to turn the photos/screenshots of class schedules into a digital form, utilizing the “Master Busy Map” logic \cite{paper_7} to further develop an algorithm to discover the common available of many class schedules.

\hspace{1.6 cm}Course feedback is only shared when students ask questions in LINE group chats and there is no platform where information can be accessed any time. Several systems such as online course platforms and academic planning tools have been developed to address individual aspects of this problem. For example, “RateMyProfessors” website \cite{ratemyprofessors}, developed to allow students to share their opinions and comments of their professors with colleagues.

\hspace{1.6 cm}Students also face difficulties when planning their academic paths. Most students rely on the student handbook to check course requirements, prerequisites, and credit conditions. However, the handbook only serves as a reference document and does not provide interactive guidance. As a result, many students become confused about which general education courses they still need to take, which electives are required, and how their course selections affect their graduation timeline. It is common for students to enroll in available courses without a clear long-term plan, which may lead to delayed graduation or inefficient course sequencing. Stellic and DegreesWork \cite{stellic} \cite{degreeworks}, academic planning websites, designed to let students stay informed of their academic performance and possible academic career path that they can choose.  However, most of these systems focus on single functions and lack certain capability to solve our addressed issues, for example, \cite{ratemyprofessors} lack a lot of information about EGCI curriculum structures, Stellic and DegreesWork \cite{stellic} \cite{degreeworks} use text-based academic planning, which \cite{trippel2025developing} addressed that graph-based representations of course prerequisites can improve student’s understanding of course syllabus.

\hspace{1.6 cm}The objective of this project is to develop EGCEye, an integrated academic planning and course rating web platform specifically for EGCI students. The system will include three main functions: a class schedule matcher using image processing techniques, a rate-my-course feature for student reviews and feedback, and a super visualized planner that guides students through their academic progress toward graduation.



\section{Objectives}
\begin{itemize}
	\item To analyze the user behaviors regarding schedule sharing, existing course review , and course planning system
	\item To design a web interface and system architecture that integrates:
	\begin{enumerate}[label=\alph*.]
	\item A Schedule Matcher interface for image uploads
	\item A Rate-My-Course database scheme for multi-section feedback
	\item A Super Planner visualization engine for prerequisite tracking on an interactive drag-and-drop interface.
	\end{enumerate}
	\item To develop a website for EGCI students to:
	\begin{enumerate}[label=\alph*.]
	\item Share their class schedules with friends to find common available times
	\item Rate courses and leave comments on the courses offered by MUIC and the EGCI faculty
	\item Create an interactive academic plan to track both courses they intend to enroll in and those they have already completed throughout their studies in the EGCI faculty
	\end{enumerate}
	\item To evaluate the website’s performance in terms of image extraction accuracy, user engagement across all functions, and the overall user experience (UX).
	
\end{itemize}

\section{Scope}
\begin{itemize}
	\item Platform Availability: The system will be developed as a web application, available specifically for EGCI students.
	\item Technology Stack
	\begin{enumerate}[label=\alph*.]
	\item Full-Stack Framework: MERN, MongoDB for database, Express.js as a backend framework, React as a frontend library, and Node.js a runtime environment to run JavaScript code
	\item Image Processing Engine: A JavaScript-based service utilizing Tesseract.js for optical character recognition and OpenCV.js for color segmentation and coordinate mapping in schedule extraction, integrated directly within the Node.js backend
	\item Programming Languages: JavaScript (ES6+) for both web architecture and image data processing
	\end{enumerate}
    
\end{itemize}

\section{Data restriction for each function}
\begin{itemize}
	\item Schedule matcher
	\begin{enumerate}[label=\alph*.]
		\item The system accepts image uploads (JPEG or PNG) of class schedules exported from Sky+  application.
		\item Date and time labels are extracted from the schedule grid by the Optical Character Recognition via Tesseract.js.
		\item Extracted schedule data is stored temporarily for the duration of the matching session, and is not retained after that.
	\end{enumerate}
	\item Rate-My-Course
	\begin{enumerate}[label=\alph*.]
		\item List of courses available to review covers Major Required Courses, Engineering Core Courses, and Computer Engineering Elective Courses offered in MUIC catalog 2022.
		\item EGCI students are allowed to provide 0 to 5 star ratings, and  leave a comment whether anonymously or not.
		\item Any available courses will also be separated by professors in case the course has more than one section, which are taught by different professors.
		\item Users may submit courses additional to the initial setting.
	\end{enumerate}
	\item Super planner
	\begin{enumerate}[label=\alph*.]
		\item List of courses available for drag-and-drop includes MUIC General Education courses, Computer Engineering Major Required Courses, Engineering Core Courses, and Computer Engineering Elective Courses in the catalog.
		\item This acts as a tool for visualization, not connected to MUIC official course enrollment.
	\end{enumerate}
	
\end{itemize}

\section{Expected Results}
\begin{enumerate}
	\item A set of requirements as a foundation for the system
	\item A user-centric interface design
	\item Web Platform, the primary deliverable integrated with three core functions
	\item Performance Evaluation based on following metrics:
	\begin{itemize}
		\item The proposed system is expected to achieve a user base of 50 or more.
		\item The Rate-My-Course reaches the number for activities: 15 for adding star-rating, and 5 for leaving comments.
		\item The Class Schedule Matcher targets a number of 15-20 screenshots uploaded to find common free times, with all being accurately matched.
		\item The Super Planner aims to facilitate at least 5-10 unique student plans.
		\item The project aims for an average User Experience (UX) rating of 4.0-4.5 out of 5 stars based on quality feedback collected from a random group of users.
	\end{itemize}
	
	
\end{enumerate}
\section{Timeline}
\newcolumntype{L}[1]{>{\raggedright\let\newline\\\arraybackslash\hspace{0pt}}m{#1}}
\newcolumntype{C}[1]{>{\centering\let\newline\\\arraybackslash\hspace{0pt}}m{#1}}
\newcolumntype{R}[1]{>{\raggedleft\let\newline\\\arraybackslash\hspace{0pt}}m{#1}}	

\begin{table}[!ht]
	\footnotesize
	\sloppy
	\centering
	\caption{Project Timeline}
	\label{tab: your-table}
	\resizebox{\textwidth}{!}{ %for cross-reference
		\begin{tabular}{|l|*{11}{c|}}
			\hline
			\multirow{2}{*}{\textbf{Plan}} & \multicolumn{11}{c|}{\textbf{Timeline 2026}} \\ \cline{2-12} &  \textbf{Feb} & \textbf{Mar} & \textbf{Apr} & \textbf{May} & \textbf{Jun} & \textbf{Jul} & \textbf{Aug} & \textbf{Sep} & \textbf{Oct} & \textbf{Nov} & \textbf{Dec} \\ \hline
			1: Literature Review & 
			\cellcolor[HTML]{000000} & & & & & & & & & &  \\ \hline
			2: System Arch. Design & 
			\cellcolor[HTML]{000000} & 
			\cellcolor[HTML]{000000} & & & & & & & & &\\ \hline
			3: UI/UX Deign &  & 
			\cellcolor[HTML]{000000} &
			\cellcolor[HTML]{000000} & & & & & & &  &\\ \hline
			4: Rate My Course & & & 
			\cellcolor[HTML]{000000} & 
			\cellcolor[HTML]{000000} & & & & & & &\\ \hline
			5: Schedule Matcher& & & & & 
			\cellcolor[HTML]{000000} &
			\cellcolor[HTML]{000000} &&&&& \\ \hline
			6: Super Planner& & &  & & & 
			\cellcolor[HTML]{000000} & 
			\cellcolor[HTML]{000000} & &&&\\ \hline
			7: Testing & & & & & & & 
			\cellcolor[HTML]{000000} & 
			\cellcolor[HTML]{000000} &
			\cellcolor[HTML]{000000} & & \\ \hline
			8: Documentation & & & & & & & & & &  \cellcolor[HTML]{000000} &
			\cellcolor[HTML]{000000}   \\ \hline
			9: Final Presentation &&&&&&&&&&& \cellcolor[HTML]{000000} \\ \hline
			
	\end{tabular}}
\end{table}